\documentclass[11pt]{article}
\usepackage[a4paper,margin=22mm]{geometry}
\usepackage{fontspec}
\setmainfont{lmroman10-regular.otf}[BoldFont=lmroman10-bold.otf,ItalicFont=lmroman10-italic.otf,BoldItalicFont=lmroman10-bolditalic.otf]
\usepackage{amsmath,amssymb,amsthm,mathtools,graphicx,xcolor,hyperref,xurl}
\setlength{\emergencystretch}{4em}
\newtheorem{theo}{Theorem}[section]
\newtheorem{lemm}[theo]{Lemma}
\newtheorem{prop}[theo]{Proposition}
\newtheorem{coro}[theo]{Corollary}
\newtheorem{corol}[theo]{Corollary}
\newtheorem{conj}[theo]{Conjecture}
\newtheorem{Proposition}[theo]{Proposition}
\theoremstyle{definition}
\newtheorem{defi}[theo]{Definition}
\newtheorem{exam}[theo]{Example}
\newtheorem{exem}[theo]{Example}
\newtheorem{rema}[theo]{Remark}
\newtheorem*{rema*}{Remark}
\newtheorem*{defi*}{Definition}
\newtheorem*{theo*}{Theorem}
\newtheorem*{lemm*}{Lemma}
\newtheorem*{prop*}{Proposition}
\newcommand{\enoncename}{}
\newtheorem{corpusenonce}[theo]{\enoncename}
\newenvironment{enonce}[1]{\renewcommand{\enoncename}{#1}\begin{corpusenonce}}{\end{corpusenonce}}
\newenvironment{enonce*}[1]{\par\medskip\noindent\textbf{#1.}\itshape}{\par\medskip}
\newenvironment{proofc}{\begin{proof}}{\end{proof}}
\newcommand{\Mc}{\mathcal{M}}

\numberwithin{equation}{section}
\def\endappendix{}
%~Mouliné par MaN_auto v.0.21.8 2019-11-25 13:52:12


\newfontfamily\CorpusDoubleStroke{Noto Sans Math}
\newcommand{\mathds}[1]{\mathord{\text{\CorpusDoubleStroke\char"1D7D9}}}
%\usepackage{subcaption}







\newcommand{\N}{\mathbb{N}}
\newcommand{\R}{\mathbb{R}}
\newcommand{\Z}{\mathbb{Z}}
\newcommand{\Q}{\mathbb{Q}}
\newcommand{\D}{\mathbb{D}}
\newcommand{\Pro}{\mathbb{P}}
\newcommand{\prob}{\mathbb{P}}
\newcommand{\E}{\mathbb{E}}
\newcommand{\bbT}{\mathbb{T}}

\DeclareMathOperator{\cross}{Cross}
\DeclareMathOperator{\cov}{Cov}
\DeclareMathOperator{\var}{Var}
\DeclareMathOperator{\Piv}{Piv}
\DeclareMathOperator{\arm}{Arm}
\DeclareMathOperator{\Circ}{Circ}
\DeclareMathOperator{\Infl}{Infl}
\DeclareMathOperator{\ann}{Ann}
\DeclareMathOperator{\fold}{Fold}
\DeclareMathOperator{\fivecross}{MultiCross}

\newcommand{\vp}{\overrightarrow{p}}
\renewcommand{\Z}{\mathbb{Z}}
\renewcommand{\N}{\mathbb{N}}
\newcommand{\un}{\mathds{1}}
\newcommand{\cond}{\, \Big| \,}

\newcommand*\calD{\mathcal{D}}
\newcommand*\calE{\mathcal{E}}
\newcommand*\calG{\mathcal{G}}
\newcommand*\calN{\mathcal{N}}
\newcommand*\calT{\mathcal{T}}
\newcommand*\calV{\mathcal{V}}
\newcommand*\sfF{\mathsf{F}}

\newcommand*\eps {\varepsilon}

\newcommand*{\dd}{\mathrm{d}}
\newcommand*{\dt}{\dd t}
\newcommand*{\dx}{\dd x}
\newcommand*{\drho}{\dd \rho}
\newcommand*{\dxi}{\dd \xi}

\graphicspath{{./figures/}}

\newcommand*{\mk}{\mkern -1mu}
\newcommand*{\Mk}{\mkern -2mu}
\newcommand*{\mK}{\mkern 1mu}
\newcommand*{\MK}{\mkern 2mu}

\hypersetup{urlcolor=purple, linkcolor=blue, citecolor=red}

\newcommand*{\romanenumi}{\renewcommand*{\theenumi}{\roman{enumi}}}
\newcommand*{\Romanenumi}{\renewcommand*{\theenumi}{\Roman{enumi}}}
\newcommand*{\alphenumi}{\renewcommand*{\theenumi}{\alph{enumi}}}
\newcommand*{\Alphenumi}{\renewcommand*{\theenumi}{\Alph{enumi}}}
\let\oldexists\exists
\renewcommand*{\exists}{\mathrel{\oldexists}}
\let\oldforall\forall
\renewcommand*{\forall}{\mathrel{\oldforall}}


\begin{document}
\section*{\raggedright 2595: A KKL influence bound for monotone Borel events under nondegenerate correlated Gaussian measures}
\noindent\textbf{Author:} Alejandro Rivera; Hugo Vanneuville\par\noindent\textbf{Source:} The critical threshold for Bargmann–Fock percolation. Annales Henri Lebesgue 3 (2020), publisher version, DOI 10.5802/ahl.29. \url{https://ahl.centre-mersenne.org/item/AHL_2020__3__169_0/}
\par\noindent\textbf{License:} CC BY 4.0, \url{https://creativecommons.org/licenses/by/4.0/}. The original publisher PDF has the explicit license notice. See the saved source/license records.
\par\noindent\textbf{Editorial material note (not author text):} The selected problem is Theorem 2.19 for monotonic Borel events under nondegenerate correlated Gaussian measures. Exact author influence and monotonicity definitions, square-root norm footnote, complete Section 5 transfer/sublinearity/boundary-section proofs, and the precisely named product Gaussian prerequisite are retained. Unrelated percolation, Fourier and discretization results are not selected. All proof text and formulas below are editable author TeX. Mathematical wording, language and proof order are retained. Only the article class, title matter and bibliography presentation are adapted for independent compilation; no proof has been generated or repaired.
\par\noindent\textbf{Editorial difficulty:} H1: The authors transfer the product-Gaussian Keller–Mossel–Sen sharp-threshold theorem to correlated measures through a covariance square root and prove the needed directional-influence sublinearity by monotone boundary sections and Gaussian density domination. This is a substantive research sharp-threshold result beyond ordinary qualifying material.
\medskip\hrule\medskip
\noindent\textbf{Author statement, complete proof and local supporting text follow in source order.}
\setcounter{section}{1}\setcounter{theo}{8}
\begin{rema}\label{r.pos_def}
In all the paper, we are going to deal with positive definite functions and matrices. The convention seems to be that, on the one hand positive definite matrices are invertible whereas semi-positive definite matrices are not necessarily. On the other hand, a function $g$ is said strictly positive definite if for any $n \in \Z_{>0}$ and any $x_1,\dots, x_n$, $(g(x_i-x_j))_{1 \leq i,j \leq n}$ is a positive definite matrix while $g$ is said positive definite if the matrices $(g(x_i-x_j))_{1 \leq i,j \leq n}$ are semi-positive definite. We will follow these conventions and hope this remark will clear up any ambiguities.
\end{rema}
\setcounter{section}{2}\setcounter{subsection}{5}\setcounter{subsubsection}{0}\setcounter{theo}{14}
\subsubsection{The FKG inequality for Gaussian vectors}

The FKG inequality is a crucial tool to apply percolation arguments. We say that a Borel subset $A\subseteq\R^n$ is \emph{increasing} if for each $x \in A$ and $y\in\R^n$ such that $y_i \geq x_i$ for every $i \in \lbrace 1,\dots, n \rbrace$, we have $y \in A$. We say that $A$ is decreasing if the complement of $A$ is increasing. The FKG inequality for Gaussian vectors was proved by Pitt and can be stated as follows.

\begin{theo}[\cite{pitt1982positively}]\label{t.FKG}
Let $X = (X_1,\dots, X_n)$ be a $n$-dimensional Gaussian vector whose correlation matrix has non-negative entries. Then, for every increasing Borel subsets $A,B \subseteq \R^n$, we have:
\[
\Pro \left[X \in A \cap B \right] \geq \Pro \left[X \in A \right] \cdot \Pro \left[X \in B \right] \,.
\]
\end{theo}

This result is the reason why we work with positively correlated Gaussian fields. Indeed, the FKG inequality is a crucial ingredient in the proof of RSW-type results. Note however that the very recent~\cite{bg_17} proves a box-crossing property without this inequality, albeit only in a discrete setting.

\subsubsection{A differential formula}\label{sss.russo}

As in the case of Bernoulli percolation, we need to introduce a notion of influence. This notion of influence is inspired by the \emph{geometric influences} studied by Keller, Mossel and Sen in the case of product spaces~\cite{keller2012geometric,keller2014geometric}. (For more about the relations between the geometric influences and the influences of Definition~\ref{d.infl} below, which are roughly the same in the case of product measures, see Subsection~\ref{ss.KMSnonproduct}.)

\begin{defi}\label{d.infl}
Let $\mu$ be a finite Borel measure on $\R^n$, let $v \in \R^n$, and let $A$ be a Borel subset of $\R^n$. The influence of $v$ on $A$ under $\mu$ is:
\[
I_{v,\mu}(A) := \underset{r \downarrow 0}{\liminf} \, \frac{\mu \left(A + [-r,r] \, v \right) - \mu \left(A \right)}{r} \in [0,+\infty] \,.
\]
Write $\left(e_1,\dots, e_n \right)$ for the canonical basis of $\R^n$. We will use the following simplified notations:
\[
I_{i,\mu}(A) := I_{e_i,\mu}(A) \,.
\]
\end{defi}
\setcounter{section}{2}\setcounter{subsection}{5}\setcounter{subsubsection}{2}\setcounter{theo}{18}
\subsubsection{A KKL-KMS (Kahn--Kalai--Linial --- Keller--Mossel--Sen) theorem for non-product Gaussian vectors}\label{sss.KMS}

One of the contributions of this paper is the derivation of a KKL theorem for non-product Gaussian vectors, namely Theorem~\ref{t.KMSnonproduct}, based on a similar result for product Gaussian vectors by Keller, Mossel and Sen~\cite{keller2012geometric}. (This similar result by~\cite{keller2012geometric} is stated in Theorem~\ref{t.KMS} of our paper\footnote{As we will explain, Theorem~\ref{t.KMS} is a simple consequence of Item~$(1)$ of Theorem~1.5 of~\cite{keller2012geometric}.}.) Actually, with our techniques of Section~\ref{s.KMS}, most of the results from~\cite{keller2012geometric} could be extended to non-product Gaussian vectors (and to monotonic events).

\begin{theo}\label{t.KMSnonproduct}
There exists an absolute constant $c > 0$ such that the following holds: Let $n \in \Z_{>0}$, let $\Sigma$ be a $n \times n$ symmetric positive definite matrix\footnote{Remember Remark~\ref{r.pos_def}: this means in particular that $\Sigma$ is non-degenerate.} and let $\mu = \mathcal{N}(0,\Sigma)$. Also, let $\sqrt{\Sigma}$ be a symmetric square root of $\Sigma$ and write $\| \sqrt{\Sigma} \|_{\infty,op}$ for the operator norm of $\sqrt{\Sigma}$ for the infinite norm\footnote{I.e. $\| \sqrt{\Sigma} \|_{\infty,op} \Mk=\Mk \sup_{x\in \R^n \setminus \{0\}} \frac{|\sqrt{\Sigma}(x) |_\infty}{| x |_\infty}$. Equivalently, $\| \sqrt{\Sigma} \|_{\infty,op} \Mk=\Mk\max_{i \in \lbrace 1,\dots, n \rbrace} \sum_{j=1}^n |\sqrt{\Sigma}(i,j)|$.} on $\R^n$. For every $A$ monotonic Borel subset of $\R^n$ we have:
\[
\sum_{i=1}^n I_{i,\mu}(A) \geq c \, \| \sqrt{\Sigma} \|_{\infty,op}^{-1} \, \mu(A) \, (1-\mu(A)) \, \sqrt{\log_+ \left(\frac{1}{\| \sqrt{\Sigma} \|_{\infty,op} \cdot \max_{i \in \lbrace 1,\dots, n \rbrace} I_{i,\mu}(A)} \right)} \,.
\]
\end{theo}

This theorem holds for a wider class of sets. For instance, it holds for semi-algebraic\footnote{We say that a set $A\subset\R^n$ is semi-algebraic if it belongs to the Boolean algebra generated by sets of the form $P^{-1}\left(]0,+\infty[\right)$ where $P\in\R[X_1,\dots,X_n]$.}, see Appendix A of~\cite[Chapter~7]{alelele} or Appendix A of~\cite[Chapter~2]{hugogogo}. Since we need it only for monotonic events, we did not try to identify the weakest possible assumptions for the property to hold. In particular, we have not found any examples of Borel sets for which the theorem does not hold.

A way to understand the constant $\| \sqrt{\Sigma} \|_{\infty,op}$ is to see what happens in the extreme cases:
\begin{itemize}
\item If $\Sigma$ is the identity matrix, then $\| \sqrt{\Sigma} \|_{\infty,op}=1$ and the above result corresponds to the product case from~\cite{keller2012geometric}.
\item If $\Sigma_{i,j}=1$ for every $i,j$ (which corresponds to the fact that, if $X \sim \Sigma$, then $X_i=X_j$ for every $i,j$) then $\| \sqrt{\Sigma} \|_{\infty,op}=n$ which is coherent since there is no threshold phenomenon for a single variable.
\end{itemize}

The proof of Theorem~\ref{t.KMSnonproduct} is organized as follows: In Subsection~\ref{ss.KMSnonproduct}, we explain how to deduce Theorem~\ref{t.KMSnonproduct} from a sub-linear property of the influences (see Proposition~\ref{p.sublin}) and from the results of~\cite{keller2012geometric}. In Subsection~\ref{ss.monotonic}, we prove the sub-linearity property for monotonic subsets.
\setcounter{section}{4}
\section{A KKL theorem for biased Gaussian vectors: the proof of Theorem~\ref{t.KMSnonproduct}}\label{s.KMS}

In this section we prove Theorem~\ref{t.KMSnonproduct}. The proofs presented here do not rely on any results of the other sections. Recall that we use the following definition for influence. Given a vector $v\in\R^n$ and a Borel probability measure $\mu$ on $\R^n$ the influence of $v$ on $A$ under $\mu$ is
\[
I_{v,\mu}(A)=\underset{r \downarrow 0}{\liminf}\frac{\mu \left(A+[-r,r]v \right) - \mu \left(A \right)}{r} \in [0,+\infty]\,.
\]
\subsection{Sub-linearity of influences implies Theorem~\ref{t.KMSnonproduct}}\label{ss.KMSnonproduct}

The aim of this subsection is to prove Theorem~\ref{t.KMSnonproduct}. That this theorem holds for product Gaussian measures is a result by Keller, Mossel and Sen~\cite{keller2012geometric} (see Corollary~\ref{c.KMS} below). In order to extend this result to general Gaussian measures, we use a sub-linearity property for influences. In the present subsection, we state this property and use it to derive Theorem~\ref{t.KMSnonproduct} from the product case. In Subsection~\ref{ss.monotonic}, we prove the sub-linearity property for monotonic sets.

\paragraph{A KKL theorem for product Gaussian vectors} The authors of~\cite{keller2012geometric} introduce and study the notion of \emph{geometric influences} which are defined as follows: Let $\nu$ be a probability measure on $\R$ and let $A$ be a Borel subset of $\R^n$. If $i \in \lbrace 1,\dots, n \rbrace$, let:
\[
A_i^x := \lbrace y \in \R \, : \, (x_1,\dots, x_{i-1}, y, x_{i+1},\dots, x_n) \in A \rbrace \,.
\]
The geometric influence of $i$ on $A$ under the measure $\nu^{\otimes n}$ is:
\[
I^\mathcal{G}_{i,\nu}(A) := \E_{x \sim \nu^{\otimes n}} \left[\nu^+(A_i^x) \right] \in [0,+\infty] \,,
\]
where $\nu^+$ is the lower Minkowski content, defined as follows: for all $B \subseteq \R$ Borel,
\[
\nu^+(B) := \underset{r \downarrow 0}{\liminf} \frac{\nu\left(B+[-r,r]\right) - \nu \left(B\right)}{r} \in [0,+\infty] \,.
\]


In the case where $\mu=\nu^{\otimes n}$, $I^\calG_{i,\nu}$ and $I_{i,\mu}$ are closely related. Indeed, firstly, by Fubini's Theorem and Fatou's lemma, for each Borel subset $A\subset\R^n$ and each $i\in\{1,\dots,n\}$,
\begin{equation}\label{e.KMS.Fatou}
I_{i,\nu^{\otimes n}}(A)=\liminf_{r\downarrow 0} \E_{x \sim \nu^{\otimes n}} \left[\frac{\nu(A_i^x+[-r,r]) - \nu(A)}{r} \right]\geq I^\calG_{i,\nu}(A) \,.
\end{equation}

While the reverse inequality seems not true in general, we expect it to hold for a wide class of events. In particular, our Lemma~\ref{l.conv} and~(2.6) from~\cite{keller2012geometric} imply that this is the case for monotonic events. Since it is not useful to us, we do not investigate this matter any further.

We will need the following result, which is a direct consequence of Item~(1) of Theorem~1.5 of~\cite{keller2012geometric}. Several results of this type can also be found in the more recent~\cite{cordero2012hypercontractive} (see for instance the paragraph above Corollary~7 therein).

\begin{theo}\label{t.KMS}
There exists an absolute constant $c > 0$ such that the following holds:

Let $\nu = \mathcal{N}(0,1)$ and let $A$ be a Borel measurable subset of $\R^n$. Then:
\[
\sum_{i=1}^n I^\mathcal{G}_{i,\nu}(A) \geq c \, \nu^{\otimes n}(A) \, (1-\nu^{\otimes n}(A)) \, \sqrt{\log_+ \left(\frac{1}{\max_{i \in \lbrace 1,\dots, n \rbrace} I^\mathcal{G}_{i,\nu}(A)}\right) } \,.
\]
\end{theo}

\begin{coro}\label{c.KMS}
Let $\nu$ and $A$ be as in Theorem~\ref{t.KMS}. Then:
\[
\sum_{i=1}^n I_{i,\nu^{\otimes n}}(A) \geq c \, \nu^{\otimes n}(A) \, (1-\nu^{\otimes n}(A)) \, \sqrt{\log_+ \left(\frac{1}{\max_{i \in \lbrace 1,\dots, n \rbrace} I_{i,\nu^{\otimes n}}(A)}\right) } \,.
\]
\end{coro}

\begin{proof}
This is a direct consequence of Theorem~\ref{t.KMS} and Equation~\eqref{e.KMS.Fatou}.
\end{proof}

Now, let us state a sub-linearity property for the influences that we will be proved in Subsection~\ref{ss.monotonic}.

\begin{prop}\label{p.sublin}
Let $\Sigma$ be a $n \times n$ symmetric positive definite matrix, let $\mu \sim \mathcal{N}(0,\Sigma)$ and let $\left(e_1,\dots, e_n \right)$ be the canonical basis of $\R^n$. Also, let $v = \sum_{j=1}^n v_i \, e_i \in \R^n$. For every $i \in \lbrace 1,\dots, n \rbrace$ and every $A$ monotonic Borel subset of $\R^n$, we have:
\begin{equation}\label{e.sublinrec}
I_{v,\mu}(A) \leq \sum_{i=1}^n |v_i| \, I_{i,\mu}(A) \,.
\end{equation}
\end{prop}

We expect the above result to hold for a larger class of Borel subsets $A$ (in particular, we have not found any examples of Borel sets for which this does not hold) and not only for the decomposition on the canonical basis. See~\cite[Chapter~7, Appendix~A]{alelele} or~\cite[Chapter~2, Appendix~A]{hugogogo} for a proof of the inequality $I_{u+v,\mu}(A)\leq I_{u,\mu}(A)+I_{v,\mu}(A)$ for any $v,w$ when $A$ is a semi-algebraic set.

Let us now derive Theorem~\ref{t.KMSnonproduct} from Corollary~\ref{c.KMS} using Proposition\ref{p.sublin}.

\begin{proof}[Proof of Theorem~\ref{t.KMSnonproduct}] Let $\sqrt{\Sigma}$ be a symmetric square root of $\Sigma$ and let $\nu=\mathcal{N}(0,1)$. Then, $\mu=\calN(0,\Sigma)$ is the pushforward measure of $\nu^{\otimes n}$ by $\sqrt{\Sigma}$. Thanks to Corollary~\ref{c.KMS} (applied to the event $\sqrt{\Sigma}^{-1}(A)$), it is sufficient to prove the following claim.
\end{proof}

\begin{enonce}{Claim}
Let $\nu = \mathcal{N}(0,1)$. We have:
\begin{equation}\label{e.cons_sublin_1}
\max_{j\in\lbrace 1,\dots, n\rbrace} I_{j,\nu^{\otimes n}}(\sqrt{\Sigma}^{-1} \left(A \right)) \leq \| \sqrt{\Sigma} \|_{\infty,op} \cdot \max_{i \in \lbrace 1,\dots, n \rbrace} I_{i,\mu}(A) \,.
\end{equation}
Moreover:
\begin{equation}\label{e.cons_sublin_2}
\sum_{j=1}^n I_{j,\nu^{\otimes n}}(\sqrt{\Sigma}^{-1} \left(A \right)) \leq \| \sqrt{\Sigma} \|_{\infty,op} \cdot \sum_{i=1}^n I_{i,\mu}(A) \,.
\end{equation}
\end{enonce}

\begin{proof}
For each $j\in\lbrace 1,\dots, n\rbrace$:
\begin{align*}
I_{j,\nu^{\otimes n}}(\sqrt{\Sigma}^{-1} \left(A \right)) & = \underset{r \downarrow 0}{\liminf} \frac{\nu^{\otimes n} \left(\sqrt{\Sigma}^{-1} \left(A \right) + [-r,r] \, e_j \right) - \nu^{\otimes n} \left(\sqrt{\Sigma}^{-1} \left(A \right) \right)}{r}\\
& = \underset{r \downarrow 0}{\liminf} \frac{\mu \left(A + [-r,r] \, \sqrt{\Sigma} \cdot e_j \right) - \mu \left(A \right)}{r}\\
& = I_{\sqrt{\Sigma} \cdot e_j,\mu}(A) \,.
\end{align*}
By using Proposition~\ref{p.sublin}, we obtain that:
\begin{equation}\label{e.cons_sublin_inter}
I_{j,\nu^{\otimes n}}(\sqrt{\Sigma}^{-1} \left(A \right)) \leq \sum_{i = 1}^n |\sqrt{\Sigma}(i,j)| I_{i,\mu}(A) \,.
\end{equation}
Hence (since $\sqrt{\Sigma}$ is symmetric):
\begin{align*}
I_{j,\nu^{\otimes n}}(\sqrt{\Sigma}^{-1} \left(A \right))& \leq \sum_{i = 1}^n |\sqrt{\Sigma}(i,j)| \max{i \in \lbrace 1,\dots, n \rbrace} I_{i,\mu}(A) \\
&\leq \|\sqrt{\Sigma} \|_{\infty,op} \, \max_{i \in \lbrace 1,\dots, n \rbrace} I_{i,\mu}(A) \,.
\end{align*}
We obtain~\eqref{e.cons_sublin_1} by taking the supremum over $j$. Inequality~\eqref{e.cons_sublin_inter} also implies that:
\begin{align*}
\sum_{j=1}^n I_{\sqrt{\Sigma} \cdot e_j,\mu}(A) & \leq \sum_{j=1}^n \sum_{i=1}^n |\sqrt{\Sigma}(i,j)| \, I_{i,\mu}(A)
 = \sum_{i=1}^n I_{i,\mu}(A) \sum_{j=1}^n |\sqrt{\Sigma}(i,j)|\\
& \leq \|\sqrt{\Sigma}\|_{\infty,op} \cdot \sum_{i=1}^n I_{i,\mu}(A) \,,
\end{align*}
which is~\eqref{e.cons_sublin_2}.
\end{proof}

\subsection{Sub-linearity of influences for monotonic events}\label{ss.monotonic}

The proof of Proposition~\ref{p.sublin} relies on the following lemma.

\begin{lemm}\label{l.conv}
Let $\Sigma$ be a $n \times n$ symmetric positive definite matrix and let $\mu = \mathcal{N}(0,\Sigma)$. Moreover, let $v \in \R^n$. For every $A$ monotonic Borel subset of $\R^n$ and every $i \in \lbrace 1,\dots, n \rbrace$, we have:
\[
\exists \lim_{r \downarrow 0} \frac{\mu \left(A + [-r,r] e_i + [-r,r]v \right) - \mu \left(A + [-r,r]v \right)}{r} = I_{i,\mu}(A) \,.
\]
\end{lemm}

We first prove Proposition~\ref{p.sublin} using Lemma~\ref{l.conv}.

\begin{proof}[Proof of Proposition~\ref{p.sublin}]
Let $A$ be a monotonic Borel subset of $\R^n$ and fix $v \in \R^n$. Let $v^i=\sum_{k=i}^n v_k e_k$. We will prove that, for every $i \in \{ 1,\dots, n-1 \}$:
\begin{equation}\label{e.sublin}
I_{v^i,\mu}(A) \leq |v_i| \, I_{i,\mu}(A) + I_{v^{i+1},\mu}(A) \,.
\end{equation}
The result will then follow directly by induction (and since $I_{v^{n},\mu}(A)= I_{n,\mu}(A)$). For any $w_1,w_2 \subseteq \R^n$, we have $[-r,r](w_1+w_2) \subseteq [-r,r]w_1+[-r,r]w_2$. Hence:
\begin{align*}
\mu \left(A + [-r,r]v^i \right) &=\mu \left(A + [-r,r] \left(v_i e_i + v^{i+1} \right) \right)\\
&\leq \mu \left(A + [-|v_i| r,|v_i| r] e_i + [-r,r] v^{i+1} \right)\\
&= \mu \left(A + [-|v_i| r,|v_i| r] e_i + [-r,r] v^{i+1} \right)\\
&\qquad - \mu \left(A + [-r,r] v^{i+1} \right) + \mu \left(A + [-r,r] v^{i+1} \right) \,.
\end{align*}

By Lemma~\ref{l.conv} we have:
\[
\frac{\mu \left(A + [-|v_i| r,|v_i| r] \, e_i + [-r,r] \, v^{i+1} \right) - \mu \left(A + [-r,r]v^{i+1} \right)}{r} \underset{r \downarrow 0}{\longrightarrow} |v_i| \, I_{i,\mu}(A) \,.
\]
Equation~\eqref{e.sublin} follows.
\end{proof}

\begin{proof}[Proof of Lemma~\ref{l.conv}] We are inspired by the proof of Proposition~1.3 in~\cite{keller2012geometric}. We write the proof for $A$ decreasing since the proof for $A$ increasing is identical. Also, we prove the result in the case $i = n$. We write $\tilde{x}$ for the first $(n-1)$ coordinates of any $x\in\R^n$. Let $A_r = A + [-r,r]v$, note that $A_r$ is decreasing, and write for any $\tilde{x}\in\R^{n-1}$:
\begin{align*}
s(\tilde{x}) &:= \sup \lbrace x_n \in \R \, : \, (\tilde{x}, x_n) \in A \rbrace \in [-\infty,+\infty] \,,\\
s_r(\tilde{x}) &:= \sup \lbrace x_n \in \R \, : \, (\tilde{x}, x_n) \in A_r \rbrace \in [-\infty,+\infty] \,.
\end{align*}
Let $\lambda$ be the density function of $\mu$. Since $A$ and $A_r$ are decreasing, we have:
\begin{align}\label{e.mon.integral}
\frac{\mu (A + [-r,r] e_n) - \mu (A)}{r} &= \frac{1}{r} \int_{\R^{n-1}} \left(\int_{s(\tilde{x})}^{s(\tilde{x})+r} \lambda(\tilde{x},x_n) \dd x_n\right)\dd\tilde{x} \,, \nonumber\\
\frac{\mu (A_r + [-r,r] e_n) - \mu (A_r)}{r} &= \frac{1}{r} \int_{\R^{n-1}} \left(\int_{s_r(\tilde{x})}^{s_r(\tilde{x})+r} \lambda(\tilde{x},x_n) \dd x_n\right)\dd\tilde{x} \,,
\end{align}
where by convention $\int_{-\infty}^{-\infty+r} = \int_{+\infty}^{+\infty+r} = 0$. For each $\tilde{x}\in\R^{n-1}$ let:
\[
g_1(\tilde{x})=\sup_{x_n\in\R}\lambda(\tilde{x},x_n);\: \: g_2(\tilde{x})=\sup_{x_n\in\R}\left|\frac{\partial\lambda}{\partial x_n}(\tilde{x},x_n)\right|.
\]
Direct computation shows that $g_1,g_2\in L^1(\R^{n-1})$. By the mean value inequality, for each $\tilde{x}\in\R^{n-1}$:
\[
\left|\frac{1}{r}\left(\int_{s(\tilde{x})}^{s(\tilde{x})+r} \lambda(\tilde{x},x_n) \dx_n\right)-\lambda(\tilde{x},s(\tilde{x}))\right|+\left|\frac{1}{r}\left(\int_{s_r(\tilde{x})}^{s_r(\tilde{x})+r} \lambda(\tilde{x},x_n) \dd x_n\right)-\lambda(\tilde{x},s_r(\tilde{x}))\right|
\]
is no greater than $2rg_2(\tilde{x})$. Combining this with Equation~\eqref{e.mon.integral} we get:
\begin{multline*}
\left|\frac{\mu (A + [-r,r] e_n) - \mu (A)}{r}- \frac{\mu (A_r + [-r,r] e_n) - \mu (A_r)}{r} \right|\\
\leq \left|\int_{\R^{n-1}}\lambda(\tilde{x},s(\tilde{x})) -\lambda(\tilde{x},s_r(\tilde{x})) \, \dd\tilde{x}\right|+2r\int_{\R^{n-1}}g_2(\tilde{x}) \, \dd\tilde{x}\,.
\end{multline*}
Since $g_2\in L^1(\R^{n-1})$, the second integral in the last inequality is finite and independent of $r$. Moreover:
\[
\int_{\R^{n-1}}\lambda(\tilde{x},s_r(\tilde{x})) -\lambda(\tilde{x},s(\tilde{x})) \, \dd\tilde{x}=\int_{\R^{n-1}}\int_{s(\tilde{x})}^{s_r(\tilde{x})}\frac{\partial\lambda}{\partial x_n}(\tilde{x},x_n) \, \dd x_n \dd\tilde{x}\,.
\]
Since $\frac{\partial\lambda}{\partial x_n}\in L^1(\R^n)$, by dominated convergence, all that remains is to show that for a.e. $\tilde{x}\in\R^{n-1}$: $s_r(\tilde{x}) \longrightarrow_{r\downarrow 0} s(\tilde{x})$. Since for each $\tilde{x}$ the sequence $s_r(\tilde{x})$ is decreasing, it converges to some $s_\infty(\tilde{x}) \geq s(\tilde{x})$. Let us prove that, for a.e. $\tilde{x}\in\R^{n-1}$, $s_\infty(\tilde{x})=s(\tilde{x})$. To do so, first note that, since $A$ is decreasing, we have:
\[
0 \leq \mu(A_r) - \mu(A) \leq \Pro \left[X - \sum_{i=1}^n r |v_i|e_i \in A \right] - \Pro \left[X \in A \right] \,,
\]
where $X \sim \calN(0,\Sigma)$. By dominated convergence, the right hand side tends to $0$ when $r\rightarrow 0$. Now, note that:
\[
\mu(A_r)-\mu(A) = \int_{\R^{n-1}} \int_{s(\tilde{x})}^{s_r(\tilde{x})} \lambda(\tilde{x},x_n) \, \dd x_n \dd\tilde{x} \,.
\]
Hence, by Fatou's lemma:
\[
0 = \int_{\R^{n-1}} \int_{s(\tilde{x})}^{s_\infty(\tilde{x})} \lambda(\tilde{x},x_n) \, \dd x_n \dd\tilde{x} \,.
\]
Since the $\lambda$ takes only positive values, this implies that for a.e. $\tilde{x}\in\R^{n-1}$, $s_\infty(\tilde{x})=s(\tilde{x})$.
\end{proof}
\begin{thebibliography}{MMMMMM}
\bibitem[Ale96]{alex_96} Alexander, Kenneth S. Boundedness of level lines for two-dimensional random fields., Ann. Probab., Volume 24 (1996) no. 4, pp. 1653-1674 \href{https://ahl.centre-mersenne.org/item/AHL_2020__3__169_0/#alex_96}{Publisher reference}.
\bibitem[ATT18]{ahlberg2016sharpness} Ahlberg, Daniel; Tassion, Vincent; Teixeira, Augusto Sharpness of the phase transition for continuum percolation in \ensuremath{\mathbb{R}} 2 , Probab. Theory Relat. Fields, Volume 172 (2018) no. 1-2, pp. 525-581 \href{https://ahl.centre-mersenne.org/item/AHL_2020__3__169_0/#ahlberg2016sharpness}{Publisher reference}.
\bibitem[AW09]{azws} Azaïs, Jean-Marc; Wschebor, Mario Level sets and extrema of random processes and fields, John Wiley \& Sons, 2009, xii+393 pages \href{https://ahl.centre-mersenne.org/item/AHL_2020__3__169_0/#azws}{Publisher reference}.
\bibitem[BDC12]{beffara2012self} Beffara, Vincent; Duminil-Copin, Hugo The self-dual point of the two-dimensional random-cluster model is critical for q\ensuremath{\geq}1, Probab. Theory Relat. Fields, Volume 153 (2012) no. 3, pp. 511-542 \href{https://ahl.centre-mersenne.org/item/AHL_2020__3__169_0/#beffara2012self}{Publisher reference}.
\bibitem[BG17a]{bg_16} Beffara, Vincent; Gayet, Damien Percolation of random nodal lines, Publ. Math., Inst. Hautes Étud. Sci., Volume 126 (2017) no. 1, pp. 131-176 \href{https://ahl.centre-mersenne.org/item/AHL_2020__3__169_0/#bg_16}{Publisher reference}.
\bibitem[BG17b]{bg_17} Beffara, Vincent; Gayet, Damien Percolation without FKG (2017) (\url{https://arxiv.org/abs/1710.10644}) \href{https://ahl.centre-mersenne.org/item/AHL_2020__3__169_0/#bg_17}{Publisher reference}.
\bibitem[BKK + 92]{bourgain1992influence} Bourgain, Jean; Kahn, Jeff; Kalai, Gil; Katznelson, Yitzhak; Linial, Nathan The influence of variables in product spaces, Isr. J. Math., Volume 77 (1992) no. 1-2, pp. 55-64 \href{https://ahl.centre-mersenne.org/item/AHL_2020__3__169_0/#bourgain1992influence}{Publisher reference}.
\bibitem[BM18]{bm_17} Beliaev, Dmitry; Muirhead, Stephen Discretisation schemes for level sets of planar Gaussian fields, Commun. Math. Phys., Volume 359 (2018) no. 3, pp. 869-913 \href{https://ahl.centre-mersenne.org/item/AHL_2020__3__169_0/#bm_17}{Publisher reference}.
\bibitem[BMW17]{bmw_17} Beliaev, Dmitry; Muirhead, Stephen; Wigman, Igor Russo-Seymour-Welsh estimates for the Kostlan ensemble of random polynomials (2017) (\url{https://arxiv.org/abs/1709.08961}) \href{https://ahl.centre-mersenne.org/item/AHL_2020__3__169_0/#bmw_17}{Publisher reference}.
\bibitem[BR06a]{bollobas2006critical} Bollobás, Béla; Riordan, Oliver The critical probability for random Voronoi percolation in the plane is 1/2, Probab. Theory Relat. Fields, Volume 136 (2006) no. 3, pp. 417-468 \href{https://ahl.centre-mersenne.org/item/AHL_2020__3__169_0/#bollobas2006critical}{Publisher reference}.
\bibitem[BR06b]{bollobas2006percolation} Bollobás, Béla; Riordan, Oliver Percolation, Cambridge University Press, 2006, x+323 pages \href{https://ahl.centre-mersenne.org/item/AHL_2020__3__169_0/#bollobas2006percolation}{Publisher reference}.
\bibitem[BR06c]{bollobas2006sharp} Bollobás, Béla; Riordan, Oliver Sharp thresholds and percolation in the plane, Random Struct. Algorithms, Volume 29 (2006) no. 4, pp. 524-548 \href{https://ahl.centre-mersenne.org/item/AHL_2020__3__169_0/#bollobas2006sharp}{Publisher reference}.
\bibitem[BR06d]{bollobas2006short} Bollobás, Béla; Riordan, Oliver A short proof of the Harris–Kesten theorem, Bull. Lond. Math. Soc., Volume 38 (2006) no. 3, pp. 470-484 \href{https://ahl.centre-mersenne.org/item/AHL_2020__3__169_0/#bollobas2006short}{Publisher reference}.
\bibitem[BS07]{bs_07} Bogomolny, Eugene; Schmit, Charles Random wavefunctions and percolation, J. Phys. A, Math. Theor., Volume 40 (2007) no. 47, pp. 14033-14043 \href{https://ahl.centre-mersenne.org/item/AHL_2020__3__169_0/#bs_07}{Publisher reference}.
\bibitem[CEL12]{cordero2012hypercontractive} Cordero-Erausquin, Dario; Ledoux, Michel Hypercontractive measures, Talagrand’s inequality, and influences, Geometric aspects of functional analysis (Lecture Notes in Mathematics), Volume 2050, Springer, 2012, pp. 169-189 \href{https://ahl.centre-mersenne.org/item/AHL_2020__3__169_0/#cordero2012hypercontractive}{Publisher reference}.
\bibitem[CL09]{cheney2009course} Cheney, Elliott Ward; Light, William Allan A course in approximation theory, Graduate Studies in Mathematics, 101, American Mathematical Society, 2009 \href{https://ahl.centre-mersenne.org/item/AHL_2020__3__169_0/#cheney2009course}{Publisher reference}.
\bibitem[DCRT17]{duminil2017exponential} Duminil-Copin, Hugo; Raoufi, Aran; Tassion, Vincent Exponential decay of connection probabilities for subcritical Voronoi percolation in \ensuremath{\mathbb{R}} d , Probab. Theory Relat. Fields, Volume 173 (2017) no. 1-2, pp. 479-490 \href{https://ahl.centre-mersenne.org/item/AHL_2020__3__169_0/#duminil2017exponential}{Publisher reference}.
\bibitem[DCRT18]{duminil2016new} Duminil-Copin, Hugo; Raoufi, Aran; Tassion, Vincent A new computation of the critical point for the planar random cluster model with q\ensuremath{\geq}1, Ann. Inst. Henri Poincaré, Probab. Stat., Volume 54 (2018) no. 1, pp. 422-436 \href{https://ahl.centre-mersenne.org/item/AHL_2020__3__169_0/#duminil2016new}{Publisher reference}.
\bibitem[DCRT19]{duminil2017sharp} Duminil-Copin, Hugo; Raoufi, Aran; Tassion, Vincent Sharp phase transition for the random-cluster and Potts models via decision trees, Ann. Math., Volume 189 (2019) no. 1, pp. 75-99 \href{https://ahl.centre-mersenne.org/item/AHL_2020__3__169_0/#duminil2017sharp}{Publisher reference}.
\bibitem[GG06]{graham2006influence} Graham, Benjamin T.; Grimmett, Geoffrey R. Influence and sharp-threshold theorems for monotonic measures, Ann. Probab., Volume 34 (2006) no. 5, pp. 1726-1745 \href{https://ahl.centre-mersenne.org/item/AHL_2020__3__169_0/#graham2006influence}{Publisher reference}.
\bibitem[GG11]{graham2011sharp} Graham, Benjamin T.; Grimmett, Geoffrey R. Sharp thresholds for the random-cluster and Ising models, Ann. Appl. Probab., Volume 21 (2011) no. 1, pp. 240-265 \href{https://ahl.centre-mersenne.org/item/AHL_2020__3__169_0/#graham2011sharp}{Publisher reference}.
\bibitem[Gra09]{grafakos} Grafakos, Loukas Classical Fourier analysis, Graduate Texts in Mathematics, 249, Springer, 2009 \href{https://ahl.centre-mersenne.org/item/AHL_2020__3__169_0/#grafakos}{Publisher reference}.
\bibitem[Gri99]{grimmett1999percolation} Grimmett, Geoffrey R. Percolation, Grundlehren der Mathematischen Wissenschaften, 321, Springer, 1999 \href{https://ahl.centre-mersenne.org/item/AHL_2020__3__169_0/#grimmett1999percolation}{Publisher reference}.
\bibitem[Gri10]{grimmett2010probability} Grimmett, Geoffrey R. Probability on graphs. Random processes on graphs and lattices, Institute of Mathematical Statistics Textbooks, 1, Cambridge University Press, 2010 \href{https://ahl.centre-mersenne.org/item/AHL_2020__3__169_0/#grimmett2010probability}{Publisher reference}.
\bibitem[GS15]{garban2014noise} Garban, Christophe; Steif, Jeffrey Noise sensitivity of Boolean functions and percolation, Institute of Mathematical Statistics Textbooks, 5, Cambridge University Press, 2015 \href{https://ahl.centre-mersenne.org/item/AHL_2020__3__169_0/#garban2014noise}{Publisher reference}.
\bibitem[Har60]{harris1960lower} Harris, Theodore E. A lower bound for the critical probability in a certain percolation process, Proc. Camb. Philos. Soc., Volume 56 (1960) no. 01, pp. 13-20 \href{https://ahl.centre-mersenne.org/item/AHL_2020__3__169_0/#harris1960lower}{Publisher reference}.
\bibitem[Kes80]{kesten1980critical} Kesten, Harry The critical probability of bond percolation on the square lattice equals 1/2, Commun. Math. Phys., Volume 74 (1980) no. 1, pp. 41-59 \href{https://ahl.centre-mersenne.org/item/AHL_2020__3__169_0/#kesten1980critical}{Publisher reference}.
\bibitem[KKL88]{kahn1988influence} Kahn, Jeff; Kalai, Gil; Linial, Nathan The influence of variables on Boolean functions, 29th Annual Symposium on Foundations of Computer Science, IEEE (1988), pp. 68-80 \href{https://ahl.centre-mersenne.org/item/AHL_2020__3__169_0/#kahn1988influence}{Publisher reference}.
\bibitem[KMS12]{keller2012geometric} Keller, Nathan; Mossel, Elchanan; Sen, Arnab Geometric influences, Ann. Probab., Volume 40 (2012) no. 3, pp. 1135-1166 \href{https://ahl.centre-mersenne.org/item/AHL_2020__3__169_0/#keller2012geometric}{Publisher reference}.
\bibitem[KMS14]{keller2014geometric} Keller, Nathan; Mossel, Elchanan; Sen, Arnab Geometric influences II: Correlation inequalities and noise sensitivity, Ann. Inst. Henri Poincaré, Probab. Stat., Volume 50 (2014) no. 4, pp. 1121-1139 \href{https://ahl.centre-mersenne.org/item/AHL_2020__3__169_0/#keller2014geometric}{Publisher reference}.
\bibitem[MS83a]{molchanov1983percolationi} Molchanov, Stanislav A.; Stepanov, A. K. Percolation in random fields. I, Theor. Math. Phys., Volume 55 (1983) no. 2, pp. 478-484 \href{https://ahl.centre-mersenne.org/item/AHL_2020__3__169_0/#molchanov1983percolationi}{Publisher reference}.
\bibitem[MS83b]{molchanov1983percolationii} Molchanov, Stanislav A.; Stepanov, A. K. Percolation in random fields. II, Theor. Math. Phys., Volume 55 (1983) no. 3, pp. 592-599 \href{https://ahl.centre-mersenne.org/item/AHL_2020__3__169_0/#molchanov1983percolationii}{Publisher reference}.
\bibitem[MS86]{molchanov1986percolationiii} Molchanov, Stanislav A.; Stepanov, A. K. Percolation in random fields. III, Theor. Math. Phys., Volume 67 (1986) no. 2, pp. 434-439 \href{https://ahl.centre-mersenne.org/item/AHL_2020__3__169_0/#molchanov1986percolationiii}{Publisher reference}.
\bibitem[NS16]{nazarov2015asymptotic} Nazarov, Fedor; Sodin, Mikhail Asymptotic laws for the spatial distribution and the number of connected components of zero sets of Gaussian random functions, Zh. Mat. Fiz. Anal. Geom., Volume 12 (2016) no. 3, pp. 205-278 \href{https://ahl.centre-mersenne.org/item/AHL_2020__3__169_0/#nazarov2015asymptotic}{Publisher reference}.
\bibitem[Pit82]{pitt1982positively} Pitt, Loren D. Positively correlated normal variables are associated, Ann. Probab., Volume 10 (1982), pp. 496-499 \href{https://ahl.centre-mersenne.org/item/AHL_2020__3__169_0/#pitt1982positively}{Publisher reference}.
\bibitem[Riv18]{alelele} Rivera, Alejandro Mécanique statistique des champs gaussiens, Ph. D. Thesis, Univ. Grenoble Alpes (France) (2018) \href{https://ahl.centre-mersenne.org/item/AHL_2020__3__169_0/#alelele}{Publisher reference}.
\bibitem[Rod15]{rodriguez20150} Rodriguez, Pierre-François A 0-1 law for the massive Gaussian free field, Probab. Theory Relat. Fields, Volume 169 (2015) no. 3-4, pp. 901-930 \href{https://ahl.centre-mersenne.org/item/AHL_2020__3__169_0/#rodriguez20150}{Publisher reference}.
\bibitem[Rud62]{rud_fou} Rudin, Walter Fourier analysis on groups, Interscience Tracts in Pure and Applied Mathematics, 12, Interscience Publishers, 1962 \href{https://ahl.centre-mersenne.org/item/AHL_2020__3__169_0/#rud_fou}{Publisher reference}.
\bibitem[Rus78]{russo1978note} Russo, Lucio A note on percolation, Probab. Theory Relat. Fields, Volume 43 (1978) no. 1, pp. 39-48 \href{https://ahl.centre-mersenne.org/item/AHL_2020__3__169_0/#russo1978note}{Publisher reference}.
\bibitem[Rus82]{russo1982approximate} Russo, Lucio An approximate zero-one law, Probab. Theory Relat. Fields, Volume 61 (1982) no. 1, pp. 129-139 \href{https://ahl.centre-mersenne.org/item/AHL_2020__3__169_0/#russo1982approximate}{Publisher reference}.
\bibitem[RV17]{rv_rsw} Rivera, Alejandro; Vanneuville, Hugo Quasi-independence for nodal lines (2017) (\url{https://arxiv.org/abs/1711.05009,} to appear in Ann. Inst. Henri Poincaré, Probab. Stat.) \href{https://ahl.centre-mersenne.org/item/AHL_2020__3__169_0/#rv_rsw}{Publisher reference}.
\bibitem[She07]{shef_gff} Sheffield, Scott Gaussian free fields for mathematicians, Probab. Theory Relat. Fields, Volume 139 (2007) no. 3-4, pp. 521-541 \href{https://ahl.centre-mersenne.org/item/AHL_2020__3__169_0/#shef_gff}{Publisher reference}.
\bibitem[SW78]{seymour1978percolation} Seymour, Paul D.; Welsh, Dominic J. A. Percolation probabilities on the square lattice, Ann. Discrete Math., Volume 3 (1978), pp. 227-245 \href{https://ahl.centre-mersenne.org/item/AHL_2020__3__169_0/#seymour1978percolation}{Publisher reference}.
\bibitem[Tal94]{talagrand1994russo} Talagrand, Michel On Russo’s approximate zero-one law, Ann. Probab., Volume 22 (1994) no. 3, pp. 1576-1587 \href{https://ahl.centre-mersenne.org/item/AHL_2020__3__169_0/#talagrand1994russo}{Publisher reference}.
\bibitem[Tas16]{tassion2014crossing} Tassion, Vincent Crossing probabilities for Voronoi percolation, Ann. Probab., Volume 44 (2016) no. 5, pp. 3385-3398 \href{https://ahl.centre-mersenne.org/item/AHL_2020__3__169_0/#tassion2014crossing}{Publisher reference}.
\bibitem[Van18]{hugogogo} Vanneuville, Hugo Percolation dans le plan : dynamiques, pavages aléatoires et lignes nodales, Ph. D. Thesis, Univ. Lyon 1 (France) (2018) \href{https://ahl.centre-mersenne.org/item/AHL_2020__3__169_0/#hugogogo}{Publisher reference}.
\end{thebibliography}
\end{document}
